\subsection{Arithmetic completion over the proconstant base}

\begin{statement}{E.30}{Proved}{Arithmetic completion over the proconstant base kills the geometric group}
Put $U=\mathbb P^1_{\mathbb Q}-\{0,1,\infty\}$ and
$\mathcal E=\pi_1([U/S_3])$, with the positive rational tangential base point lifted to $U$.
The quotient by its geometric free subgroup $\widehat F_2$ is
$S_3\times G_{\mathbb Q}$.
For $m\geq1$, let $K_m=\mathbb Q(\mu_m)$,
$C_m=\operatorname{Gal}(K_m/\mathbb Q)$ and
$B_m=S_3\times C_m$.
There are two precise versions of the following construction:
abstract relative completion over $k=\mathbb Q$, and continuous relative completion over $k=\mathbb Q_\ell$.
Use the corresponding version throughout each formula, and let
$\mathcal A_{m,k}$ be the relative completion of $G_{\mathbb Q}\to C_m$.
The relative completion of $\mathcal E\to B_m$ is exactly
\[
 \mathcal R_{m,k}=S_3\times\mathcal A_{m,k},\qquad
 1\longrightarrow(G_{K_m})^{\rm un}_k
 \longrightarrow\mathcal A_{m,k}\longrightarrow C_m
 \longrightarrow1.                                      \tag{E.30.1}
\]
For $m\mid n$ the transition maps are induced by the identity on
$G_{\mathbb Q}$; on the displayed kernels they are induced by
$G_{K_n}\hookrightarrow G_{K_m}$.
Consequently the completion over the proconstant group scheme
$B=S_3\times\widehat{\mathbb Z}^{\times}_k$ is the specified inverse limit
\[
 \mathcal R_{B,k}=\varprojlim_{m}\mathcal R_{m,k},\qquad
 1\longrightarrow\varprojlim_m(G_{K_m})^{\rm un}_k
 \longrightarrow\mathcal R_{B,k}\longrightarrow B
 \longrightarrow1.                                      \tag{E.30.2}
\]
The natural homomorphism $\mathcal E\to\mathcal R_{B,k}(k)$ kills
the whole geometric subgroup $\widehat F_2$.
In particular its restriction to every Fermat subgroup is trivial;
it retains none of the geometric homology computed in E.14.

There is a different, explicit completion retaining that geometry.
Let $\mathcal A_\ell$ be the negatively weighted completion of
$G_{\mathbb Q}\xrightarrow{\chi_\ell}\mathbb G_m(\mathbb Q_\ell)$,
with central cocharacter $t\mapsto t^{-2}$.
For the arithmetic group $\pi_1(U)=\widehat F_2\rtimes G_{\mathbb Q}$
with the stated tangential section, its negatively weighted completion is
\[
 (\widehat F_2)^{\rm un}_{\mathbb Q_\ell}\rtimes\mathcal A_\ell.
                                                               \tag{E.30.3}
\]
The geometric factor is the free prounipotent group on two generators.
No algebraic morphism from the proconstant $B$ to $\mathbb G_m$
induces $\chi_\ell$ on its cyclotomic factor.
\end{statement}
\begin{proof}
The morphism $[U/S_3]\to[\operatorname{Spec}\mathbb Q/S_3]$
gives the stated quotient $S_3\times G_{\mathbb Q}$, since the
permutation action is defined over $\mathbb Q$. Its geometric
kernel is $\widehat F_2$. Thus the kernel $N_m$ of
$\mathcal E\to B_m$ has the tangentially split sequence
\[
 1\longrightarrow\widehat F_2\longrightarrow N_m
 \longrightarrow G_{K_m}\longrightarrow1.
\]
On $\widehat F_2^{\rm ab}=\widehat{\mathbb Z}^2$ the arithmetic
action is scalar multiplication by the full cyclotomic character.

First work continuously over $\mathbb Q_\ell$.
Let $r:N_m\to U_0(\mathbb Q_\ell)$ be any continuous homomorphism
to a unipotent algebraic group. Let $A$ be the Zariski closure of
$r(\widehat F_2)$. It is a connected unipotent normal subgroup
of the Zariski closure of $r(N_m)$; unipotent groups in characteristic
zero have no nontrivial finite algebraic quotients.
The map to $A^{\rm ab}$ factors through
$\mathbb Z_\ell^2$ and spans its Lie algebra over $\mathbb Q_\ell$.
Choose $g\in G_{K_m}$ with $\chi_\ell(g)\ne1$, possible because
the cyclotomic image of this finite-index subgroup is infinite.
Conjugation by $r(g)$ on $\operatorname{Lie}(A^{\rm ab})$
is both multiplication by $\chi_\ell(g)$ and a unipotent linear
operator. The latter assertion follows, for example, by putting
$U_0$ in upper triangular unipotent matrices and considering its
conjugation action on its Lie algebra and invariant subquotients.
It follows that $A^{\rm ab}=1$. A nonzero nilpotent Lie algebra
has nonzero abelianization, so $A=1$. Every such $r$ therefore
kills $\widehat F_2$.

For the abstract rational version, the required abelianization
is still exactly $\widehat{\mathbb Z}^2$: the abstract commutator
subgroup of a finitely generated profinite group is closed by
\cite[Theorem1.4]{dNikolovSegal2007}.
Choose $u\in\widehat{\mathbb Z}^{\times}$ with $u\equiv1\pmod m$
and $u-1=dv$, where $d\ne0$ is an integer and
$v\in\widehat{\mathbb Z}^{\times}$, as in D.3.
Explicitly use $u_p=2$ for odd $p\nmid m$,
$u_p=1+p^{v_p(m)}$ for odd $p\mid m$, and
$u_2=1+2^{\max(1,v_2(m))}$.
Cyclotomic surjectivity supplies $g\in G_{K_m}$ with $\chi(g)=u$.
For an abstract homomorphism $r:N_m\to U_0(\mathbb Q)$, there is
an equivariant surjection
\[
 \widehat{\mathbb Z}^2\otimes_{\mathbb Z}\mathbb Q
 \longrightarrow\operatorname{Lie}(A^{\rm ab}).
\]
On the source $u-1=dv$ is invertible. Therefore
$\operatorname{Ad}(r(g))-1$ on the target is surjective. It is
also nilpotent, since $r(g)$ is unipotent. The target is zero,
and the preceding nilpotent Lie algebra argument again gives $A=1$.
This proves the rational assertion without assuming continuity
of a homomorphism to $U_0(\mathbb Q)$.

For a group mapping onto a finite group $C$, its relative
completion is the extension of $C$ by the unipotent completion
of the kernel. One can construct it directly by choosing coset
representatives: multiplication on the finite disjoint union of
copies of the completed kernel is given by the induced conjugation
automorphisms and the original factor-set cocycle. This is the
construction and universal-property proof used in E.14; it applies
to the abstract and continuous categories just considered.
The result above identifies $N_m^{\rm un}$ with
$(G_{K_m})^{\rm un}$. Hence every representation defining
$\mathcal R_{m,k}$ kills $\widehat F_2$ and factors uniquely through
$S_3\times G_{\mathbb Q}$. Completing this direct product gives
$S_3\times\mathcal A_{m,k}$: the finite $S_3$ image commutes with
the arithmetic image, and these two homomorphisms have exactly
the product universal property. This proves (E.30.1) and the
specified transition maps.

An algebraic homomorphism from a proconstant profinite group
scheme to a finite-type affine group scheme factors through a
finite quotient. Indeed its finitely many coordinate generators
belong to one finite stage of the coordinate-ring direct limit;
the multiplication identities also hold at a sufficiently large
stage. Every finite-dimensional algebraic representation of $B$
therefore factors through some $B_m$. Defining relative completion
by representations whose semisimplifications are algebraic
representations of $B$ gives precisely the inverse system
(E.30.2): every such finite-dimensional representation occurs at
one of its stages. This also proves its universal property.
The transition maps on the unipotent kernels are faithfully flat:
a finite-index subgroup has Zariski-dense image in the connected
unipotent closure of the larger group's image. Inverse limits
therefore give the exact proalgebraic sequence displayed in
(E.30.2). The geometric image was already killed at every stage.

For (E.30.3), the genuine action of $G_{\mathbb Q}$ on the
geometric prounipotent group has scalar $\chi_\ell$ on its
abelianization. Its kernel on that abelianization is prounipotent,
with negative weights for $t\mapsto t^{-2}$: the degree-$d$
part of its IA Lie algebra maps a degree-one generator to a
degree-$d$ word and has weight $-2(d-1)<0$ for $d\geq2$.
Thus the universal
property of weighted completion extends this action uniquely to
$\mathcal A_\ell$; the weighted conventions and construction are
those of \cite[\S\S2--4]{E_HainMatsumoto}.
The semidirect product in (E.30.3) receives the actual arithmetic
group. Given any lift of that group into a negatively weighted
extension of $\mathbb G_m$, its restriction to $\widehat F_2$
extends uniquely by continuous unipotent completion, and its
restriction to the tangential $G_{\mathbb Q}$ extends uniquely by
weighted completion. Their equivariance holds on the dense
original images, hence algebraically. Thus these two extensions
define a unique homomorphism from the semidirect product, proving
its weighted universal property. Its geometric kernel is visibly
the free prounipotent group, proving its injection as well.
Finally an algebraic map $B\to\mathbb G_m$ has finite image by
the finite-stage argument, whereas $\chi_\ell(G_{\mathbb Q})$
is infinite. This proves the claimed obstruction to that
specific base-change comparison. We do not identify an abstract
rational completion after scalar extension with a continuous
$\ell$-adic completion; the two constructions above specify
their source categories independently.
\end{proof}
